\subsection{Spectral Theory of Operators}

% **Role:** infinite-dimensional linear algebra powered by analysis.
%
% Spectrum, resolvent, spectral radius, self-adjoint/normal operators, functional calculus preview.

For bounded operators, spectrum replaces the finite-dimensional eigenvalue list. The resolvent tracks invertibility, while the spectrum records every obstruction to it.

\begin{definition} \label{def:bounded-operator-spectrum-resolvent} \hlt{Spectrum and Resolvent}\\
Let $\mathcal{X}$ be a complex Banach space and let $T\in\LL(\mathcal{X})$.
\begin{enumerate}[label=(\roman*)]
\setlength{\itemsep}{0pt}
\item \hlt{Resolvent set}: $\rho(T)\defeq\{\lambda\in\C:\lambda I-T\text{ is bijective and }(\lambda I-T)^{-1}\in\LL(\mathcal{X})\}$.
\item \hlt{Spectrum}: $\sigma(T)\defeq\C\setminus\rho(T)$.
\end{enumerate}
\end{definition}

\begin{theorem} \label{thm:spectrum-nonempty-compact-bounded-operator} \hlt{Spectrum of a Complex Bounded Operator}\\
Let $\mathcal{X}$ be a nonzero complex Banach space and let $T\in\LL(\mathcal{X})$. Then:
\begin{enumerate}[label=(\roman*)]
\setlength{\itemsep}{0pt}
\item $\sigma(T)\neq\varnothing$;
\item $\sigma(T)$ is compact in $\C$;
\item $\sigma(T)\subseteq\{\lambda\in\C:\abs{\lambda}\leq\normop{T}\}$.
\end{enumerate}
\begin{proof} For $\lambda\in\rho(T)$, $R(\lambda)\defeq(\lambda I-T)^{-1}$. If $\mu\in\rho(T)$, $\abs{\lambda-\mu}\normop{R(\mu)}<1$; $\lambda I-T=(I+(\lambda-\mu)R(\mu))(\mu I-T)$.\\
By Corollary \ref{cor:neumann-series-banach-operator}, $\lambda\in\rho(T)$, $R(\lambda)=R(\mu)(I+(\lambda-\mu)R(\mu))^{-1}$. If $\abs{\lambda}>\normop{T}$, then $\lambda I-T=\lambda(I-\lambda^{-1}T)$.\\
By Corollary \ref{cor:neumann-series-banach-operator}, $\lambda\in\rho(T)$ and $\normop{R(\lambda)}\leq(\abs{\lambda}-\normop{T})^{-1}$.
\begin{enumerate}[label=(\roman*)]
\setlength{\itemsep}{0pt}
\item Suppose $\sigma(T)=\varnothing$. By the local Neumann formula above, $R:\C\to\LL(\mathcal{X})$ is entire.\\
For $x\in\mathcal{X}$ and $\Phi\in\mathcal{X}^*$, $f_{x,\Phi}(\lambda)\defeq\Phi(R(\lambda)x)$ is an entire scalar function.\
By the estimate above, $f_{x,\Phi}(\lambda)\to0$ as $\abs{\lambda}\to+\infty$, and continuity bounds $f_{x,\Phi}$ on compact disks.\
By Liouville's theorem, $f_{x,\Phi}=0$. By Corollary \ref{cor:natural-embedding-isometry}, bounded linear functionals separate points, so $R(0)=0$, contradicting $(-T)R(0)=I$ on the nonzero space $\mathcal{X}$.
\item By the local Neumann formula, $\rho(T)$ is open, so $\sigma(T)$ is closed in $\C$.\\
By outside-disk formula, $\sigma(T)\subseteq\{\lambda\in\C:\abs{\lambda}\leq\normop{T}\}$, hence $\sigma(T)$ is bounded. Thus $\sigma(T)$ is compact.
\item If $\abs{\lambda}>\normop{T}$, then the outside-disk condition holds. By the outside-disk Neumann argument, $\lambda\in\rho(T)$.\\
Therefore every $\lambda\in\sigma(T)$ satisfies $\abs{\lambda}\leq\normop{T}$.
\end{enumerate}
\hlt{Structural Commentary}: By Neumann, $\rho(T)$ is open and $\sigma(T)\subseteq\{\lambda:\abs\lambda\leq\normop T\}$. By Liouville, $\sigma(T)\neq\varnothing$.\\
\hlt{Key Insight}: Complex Banach spectra are nonempty compact obstructions to invertibility.
\end{proof}
\end{theorem}

\begin{definition} \label{def:spectral-radius-bounded-operator} \hlt{Spectral Radius}\\
Let $\mathcal{X}$ be a nonzero complex Banach space and let $T\in\LL(\mathcal{X})$. The \hlt{spectral radius} of $T$ is $r(T)\defeq\sup\limits_{\lambda\in\sigma(T)}\abs{\lambda}$.
\end{definition}

\begin{theorem} \label{thm:spectral-radius-formula} \hlt{Spectral Radius Formula}\\
Let $\mathcal{X}$ be a nonzero complex Banach space and let $T\in\LL(\mathcal{X})$.\\
Then the sequence $(\normop{T^n}^{1/n})_{n\geq1}$ converges in $[0,\infty)$, and $\lim\limits_{n\to\infty}\normop{T^n}^{1/n}=r(T)$.
\begin{proof}
Set $a_n\defeq\normop{T^n}$ for $n\geq0$, so $a_0=1$ and $a_{m+n}\leq a_ma_n$. If $a_m=0$ for some $m\geq1$, then $T$ is nilpotent.\\
Also $\normop{T^n}^{1/n}\to0$, $\sigma(T)=\{0\}$ by finite Neumann series and Theorem \ref{thm:spectrum-nonempty-compact-bounded-operator}.\\
Hence assume $a_n>0$ for all $n\geq1$. For fixed $m\geq1$ and $n=qm+r$ with $0\leq r<m$, $a_n^{1/n}\leq a_m^{q/n}a_r^{1/n}$.\\
By $\limsup$ in $n$, $\limsup_n a_n^{1/n}\leq a_m^{1/m}$, hence $\limsup_n a_n^{1/n}\leq\inf_{m\geq1}a_m^{1/m}\leq\liminf_n a_n^{1/n}$.\\
Thus $L\defeq\lim_n a_n^{1/n}=\inf_{m\geq1}a_m^{1/m}$. If $\abs{\lambda}>L$, choose $n$ with $a_n<\abs{\lambda}^n$.\\
Then $\normop{T^n/\lambda^n}<1$, so Corollary \ref{cor:neumann-series-banach-operator} makes $\lambda^nI-T^n$ invertible.\\
With $Q_n\defeq\sum\limits_{j=0}^{n-1}\lambda^{n-1-j}T^j$, $\lambda^nI-T^n=(\lambda I-T)Q_n=Q_n(\lambda I-T)$; as these are polynomials in $T$,\\
$(\lambda I-T)^{-1}=Q_n(\lambda^nI-T^n)^{-1}$. Thus $\lambda\in\rho(T)$, so $\sigma(T)\subseteq\{z\in\C:\abs z\leq L\}$ and $r(T)\leq L$.\\
Let $R>r(T)$, $R_0>\max\{R,\normop T\}$. For $\abs z=R_0$, by Neumann expansion, $(zI-T)^{-1}=\sum\limits_{k\geq0}z^{-k-1}T^k$.\\
By Cauchy's coefficient formula $T^n=\frac{1}{2\pi i}\int_{\abs z=R_0}z^n(zI-T)^{-1}\,dz$.\\
Since $\{z:R\leq\abs z\leq R_0\}\subseteq\rho(T)$, Cauchy's theorem moves the contour to $\abs z=R$.\\
With $M_R\defeq\sup_{\abs z=R}\normop{(zI-T)^{-1}}<+\infty$, $a_n\leq M_RR^{n+1}$, hence $L\leq R$. By letting $R\downarrow r(T)$, $L\leq r(T)$.\\
\hlt{Structural Commentary}: By submultiplicativity, $L$ exists; Neumann, $r(T)\leq L$; resolvent bounds, $L\leq r(T)$.\\
\hlt{Key Insight}: The spectral radius is asymptotic operator-power growth.
\end{proof}
\end{theorem}

\begin{definition} \label{def:rayleigh-quotient-self-adjoint-operator} \hlt{Quadratic Form and Rayleigh Quotient}\\
Let $\mathcal{H}$ be a real Hilbert space and let $T\in\LL(\mathcal{H})$ be self-adjoint.
\begin{enumerate}[label=(\roman*)]
\setlength{\itemsep}{0pt}
\item \hlt{Quadratic form}: $Q_T(h)\defeq\innerprod{Th,h}$ for $h\in\mathcal{H}$.
\item \hlt{Rayleigh quotient}: $R_T(h)\defeq\frac{\innerprod{Th,h}}{\innerprod{h,h}}$ for $h\in\mathcal{H}\setminus\{0\}$.
\end{enumerate}
A unit-sphere maximizer of $Q_T$ is equivalently a maximizer of $R_T$ on $\mathcal{H}\setminus\{0\}$.
\end{definition}

\begin{example} \label{ex:multiplication-self-adjoint-without-eigenvectors} \hlt{A Self-Adjoint Operator without Eigenvectors}\\
Let $a<b$, set $H\defeq L^2([a,b];\C)$, and define $(Tf)(x)\defeq xf(x)$.\\
With $M\defeq\max\{\abs a,\abs b\}$, $\norm{Tf}_2\leq M\norm f_2$, so $T\in\LL(H)$.\\
For $u,v\in H$, $\innerprod{Tu,v}=\int_a^b xu(x)\overline{v(x)}\,dx=\int_a^b u(x)\overline{xv(x)}\,dx=\innerprod{u,Tv}$, hence $T=T^*$.\\
If $Tf=\lambda f$ for $\lambda\in\C$, then $(x-\lambda)f(x)=0$ a.e. Let $E_\lambda\defeq\{x\in[a,b]:x=\lambda\}$.\\
Since $m(E_\lambda)=0$, $f=0$ a.e. on $[a,b]\setminus E_\lambda$ and on $E_\lambda$, hence $f=0$ in $H$.\\
Thus no nonzero $f\in H$ satisfies $Tf=\lambda f$, so $T$ has no eigenvectors.
\end{example}

\begin{lemma} \label{lem:hilbert-schmidt-compact-self-adjoint-eigenvalue} \hlt{Hilbert--Schmidt Lemma}\\
Let $0\neq\mathcal H$ be real Hilbert space, $T=T^*\in\LL(\mathcal H)$ compact, $Q_T(h)\defeq\innerprod{Th,h}$. Then $\exists\lambda\in\R$, $h_*\in\mathcal H$ such that:
\begin{enumerate}[label=(\roman*)]
\setlength{\itemsep}{0pt}
\item $\norm{h_*}=1$;
\item $Th_*=\lambda h_*$;
\item $\abs{\lambda}=\normop{T}=\sup\limits_{\norm h=1}\abs{Q_T(h)}$.
\end{enumerate}
\begin{proof} If $T=0$, choose $\norm{h_*}=1$ and set $\lambda=0$; then (i)--(iii) hold. Assume $T\neq0$ and put $\eta\defeq\normop T$.\\
By Theorem \ref{thm:self-adjoint-operator-norm-quadratic-form}, $\eta=\sup_{\norm h=1}\abs{Q_T(h)}>0$. Choose $\norm{x_n}=1$ and $s_n\in\{-1,1\}$ with $Q_{s_nT}(x_n)\to\eta$.\\
Passing to a subsequence, $s_n=s$; write $h_n\defeq x_n$ and $A\defeq sT$.\\
By Theorems \ref{thm:hilbert-weak-subsequential-compactness} and \ref{thm:compact-operator-weak-to-strong-characterization}, after passing to a further subsequence, $h_n\rightharpoonup h_*$ and $Ah_n\to Ah_*$ in norm.\\
Since $Q_A(h_n)-Q_A(h_*)=\innerprod{Ah_n-Ah_*,h_n}+\innerprod{Ah_*,h_n-h_*}\to0$, $Q_A(h_*)=\eta$. Set $\lambda\defeq s\eta$.
\begin{enumerate}[label=(\roman*)]
\setlength{\itemsep}{0pt}
\item By weak lower semicontinuity, $\norm{h_*}\leq1$.\\
Since $\abs{Q_A(g)}\leq\eta\norm g^2$ $\forall g\in\mathcal H$, $\eta=Q_A(h_*)\leq\eta\norm{h_*}^2$. As $\eta>0$, $\norm{h_*}=1$.
\item By (i), $R_A(h_*)=\eta$. For $\norm h=1$, $R_A(h)=Q_A(h)\leq\abs{Q_T(h)}\leq\eta$, so $h_*$ maximizes $R_A$.\\
For $h_0\in\mathcal H$, $\phi(t)\defeq R_A(h_*+t h_0)$ has $\phi'(0)=0$, hence $0=2\innerprod{Ah_*-\eta h_*,h_0}$.\\
Since $h_0$ is arbitrary, $Ah_*=\eta h_*$. Thus $Th_*=s\eta h_*=\lambda h_*$.
\item By the definition of $\lambda$ and by Theorem \ref{thm:self-adjoint-operator-norm-quadratic-form}, $\abs{\lambda}=\eta=\normop{T}=\sup\limits_{\norm h=1}\abs{Q_T(h)}$.
\end{enumerate}
\hlt{Structural Commentary}: By compactness, signed $Q_T$ attains $\normop T$; by stationarity, $Th_*=\lambda h_*$.\\
\hlt{Key Insight}: A compact self-adjoint norm maximizer is an eigenvector.
\end{proof}
\end{lemma}

\begin{corollary} \label{cor:compact-self-adjoint-extreme-rayleigh-eigenvalues} \hlt{Extremal Rayleigh Values Are Eigenvalues}\\
Let $\mathcal{H}$ be a nonzero real Hilbert space, let $T\in\LL(\mathcal{H})$ be compact and self-adjoint, and set $Q_T(h)\defeq\innerprod{Th,h}$, $\alpha\defeq\inf\limits_{\norm h=1}Q_T(h)$, $\beta\defeq\sup\limits_{\norm h=1}Q_T(h)$.
\begin{enumerate}[label=(\roman*)]
\setlength{\itemsep}{0pt}
    \item If $\beta>0$, then $\exists h_\beta\in\mathcal H$ such that $\norm{h_\beta}=1$ and $Th_\beta=\beta h_\beta$.
    \item If $\alpha<0$, then $\exists h_\alpha\in\mathcal H$ such that $\norm{h_\alpha}=1$ and $Th_\alpha=\alpha h_\alpha$.
    \item If $\beta>0$ and $(h_n)\subset\mathcal H$ satisfies $\norm{h_n}=1$, $Q_T(h_n)\to\beta$, then $\exists n_j\uparrow\infty$ and $h_\beta\in\mathcal H$ with $h_{n_j}\to h_\beta$ in norm, $\norm{h_\beta}=1$, and $Th_\beta=\beta h_\beta$.
\end{enumerate}
\begin{proof}
\begin{enumerate}[label=(\roman*)]
\setlength{\itemsep}{0pt}
\item Assume $\beta>0$. Choose $\norm{u_n}=1$ with $Q_T(u_n)\to\beta$.\\
By Theorems \ref{thm:hilbert-weak-subsequential-compactness} and \ref{thm:compact-operator-weak-to-strong-characterization}, after passing to a subsequence, $u_n\rightharpoonup u_*$ and $Tu_n\to Tu_*$ in norm.\\
Since $Q_T(u_n)-Q_T(u_*)=\innerprod{Tu_n-Tu_*,u_n}+\innerprod{Tu_*,u_n-u_*}\to0$, $Q_T(u_*)=\beta$.\\
By weak lower semicontinuity, $\norm{u_*}\leq1$. As $Q_T(g)\leq\beta\norm g^2$ $\forall g\in\mathcal H$, $\beta=Q_T(u_*)\leq\beta\norm{u_*}^2$; $\norm{u_*}=1$.\\
Thus $u_*$ maximizes $R_T$. For $h\in\mathcal H$, $\phi(t)\defeq R_T(u_*+t h)$ has $\phi'(0)=0$, so $0=2\innerprod{Tu_*-\beta u_*,h}$.\\
Hence $Tu_*=\beta u_*$; take $h_\beta\defeq u_*$.
\item Assume $\alpha<0$. Apply (i) to $-T$. Its upper Rayleigh value is $\beta_{-T}=\sup_{\norm h=1}\innerprod{-Th,h}=-\alpha>0$, so $\exists h_\alpha$ with $\norm{h_\alpha}=1$ and $(-T)h_\alpha=(-\alpha)h_\alpha$. Hence $Th_\alpha=\alpha h_\alpha$.
\item Assume $\beta>0$, $\norm{h_n}=1$, $Q_T(h_n)\to\beta$. By compactness step in (i), $h_{n_j}\rightharpoonup h_\beta$, $Th_{n_j}\to Th_\beta$ in norm.\\
As in (i), $Q_T(h_\beta)=\beta$, $\norm{h_\beta}=1$, and $Th_\beta=\beta h_\beta$.\\
By Theorem \ref{thm:radon-riesz-hilbert-space}, since $\norm{h_{n_j}}=\norm{h_\beta}=1$, $h_{n_j}\to h_\beta$ in norm.
\end{enumerate}
\hlt{Structural Commentary}: By compactness, attain $\beta$; by $-T$, attain $\alpha$; by Radon--Riesz, $\rightharpoonup+\norm{\cdot}\to1\Rightarrow\to$.\\
\hlt{Key Insight}: Compact self-adjoint Rayleigh extrema are eigenvalues.
\end{proof}
\end{corollary}

\begin{theorem} \label{thm:hilbert-schmidt-compact-self-adjoint-spectral-theorem} \hlt{Hilbert--Schmidt Theorem for Compact Self-Adjoint Operators}\\
Let $\mathcal{H}$ be a real Hilbert space and let $K\in\LL(\mathcal{H})$ be compact and self-adjoint.\\
If $\operatorname{rank}K=+\infty$, then $\exists$ orthonormal sequence $(\psi_k)_{k\geq1}\subset\mathcal H$ and nonzero scalars $(\lambda_k)_{k\geq1}\subset\R$ such that:
\begin{enumerate}[label=(\roman*)]
\setlength{\itemsep}{0pt}
\item $K\psi_k=\lambda_k\psi_k$ for all $k\geq1$, $\abs{\lambda_{k+1}}\leq\abs{\lambda_k}$, and $\lambda_k\to0$;
\item $\operatorname{cspan}\{\psi_k:k\geq1\}=(\ker K)^\perp$;
\item $Kh=\sum\limits_{k\geq1}\lambda_k\innerprod{h,\psi_k}\psi_k$ for every $h\in\mathcal{H}$, with norm convergence.
\end{enumerate}
If $\operatorname{rank}K<+\infty$, then:
\begin{enumerate}[label=(\roman*),resume]
\setlength{\itemsep}{0pt}
\item $\exists m\geq0$, nonzero $\lambda_1,\ldots,\lambda_m\in\R$, and an orthonormal basis $(\psi_k)_{k=1}^m$ of $\operatorname{range}K=(\ker K)^\perp$ such that $K\psi_k=\lambda_k\psi_k$ and $Kh=\sum\limits_{k=1}^m\lambda_k\innerprod{h,\psi_k}\psi_k$ for every $h\in\mathcal H$.
\end{enumerate}
\begin{proof}
\begin{enumerate}[label=(\roman*)]
\setlength{\itemsep}{0pt}
 \item Put $M_0\defeq\{0\}$ and $H_0\defeq\mathcal H$. If $K|_{H_j}\neq0$, then $H_j$ is $K$-invariant and $K|_{H_j}$ is compact self-adjoint.\\
By Lemma \ref{lem:hilbert-schmidt-compact-self-adjoint-eigenvalue}, $\psi_{j+1}\in H_j$, $\lambda_{j+1}\neq0$ with $\norm{\psi_{j+1}}=1$, $K\psi_{j+1}=\lambda_{j+1}\psi_{j+1}$, and $\abs{\lambda_{j+1}}=\normop{K|_{H_j}}$.\\
Set $M_{j+1}\defeq M_j\oplus\operatorname{span}\{\psi_{j+1}\}$, $H_{j+1}\defeq M_{j+1}^{\perp}$. Since $K(M_{j+1})\subseteq M_{j+1}$, $K=K^*$, $K(H_{j+1})\subseteq H_{j+1}$.\\
Thus the construction continues while $K|_{H_j}\neq0$, and $\abs{\lambda_{j+1}}\leq\abs{\lambda_j}$.\\
If $\operatorname{rank}K=+\infty$, it cannot stop, since $K|_{H_m}=0$ implies $\operatorname{range}K\subseteq M_m$.\\
If $\lambda_k\not\to0$, then for some $\varepsilon>0$, $\norm{K\psi_j-K\psi_k}^2=\lambda_j^2+\lambda_k^2\geq2\varepsilon^2$ for infinitely many $j\neq k$, contradiction.
\item Set $M\defeq\operatorname{cspan}\{\psi_k:k\geq1\}$. Since $K(M)\subseteq M$ and $K=K^*$, $K(M^\perp)\subseteq M^\perp$.\\
If $K|_{M^\perp}\neq0$, by Lemma \ref{lem:hilbert-schmidt-compact-self-adjoint-eigenvalue} choose $\phi\in M^\perp$, $\mu\neq0$ with $\norm{\phi}=1$ and $K\phi=\mu\phi$.\\
Then $\phi\in H_{k-1}$ $\forall k$, so $\abs{\mu}\leq\normop{K|_{H_{k-1}}}=\abs{\lambda_k}\to0$, contradiction. Hence $M^\perp\subseteq\ker K$.\\
Conversely, $x\in\ker K$ implies $\innerprod{x,\psi_k}=\lambda_k^{-1}\innerprod{Kx,\psi_k}=0$ $\forall k$, so $\ker K\subseteq M^\perp$.\\
Thus $M^\perp=\ker K$ and $M=(\ker K)^\perp$ by Theorem \ref{thm:hilbert-double-orthogonal-complement}.
\item For $h\in\mathcal H$, write $h=m+z$ with $m\in M=(\ker K)^\perp$ and $z\in\ker K$. By Theorem \ref{thm:hilbert-fourier-expansion-general},\\
$m=\sum\limits_{k\geq1}\innerprod{h,\psi_k}\psi_k$ in norm. By continuity of $K$, $Kh=Km=\sum\limits_{k\geq1}\lambda_k\innerprod{h,\psi_k}\psi_k$ in norm.
\item If $\operatorname{rank}K<+\infty$, the construction stops at some $m\geq0$; otherwise $\{\psi_k\}\subset\operatorname{range}K$ is an infinite orthonormal set in finite-dimensional $\operatorname{range}K$. Let $M_m\defeq\operatorname{span}\{\psi_1,\ldots,\psi_m\}$.\\
Since $K|_{M_m^\perp}=0$, $\operatorname{range}K\subseteq M_m$; since $\psi_k=\lambda_k^{-1}K\psi_k$, $M_m\subseteq\operatorname{range}K$. Hence $M_m=\operatorname{range}K$.\\
Also $M_m^\perp=\ker K$ by argument in (ii), so $\operatorname{range}K=(\ker K)^\perp$. For $h=m_0+z$ with $m_0\in M_m$, $z\in\ker K$,\\
by the finite orthonormal expansion of $m_0$, $Kh=Km_0=\sum\limits_{k=1}^m\lambda_k\innerprod{h,\psi_k}\psi_k$.
\end{enumerate}
\hlt{Structural Commentary}: Iterated extremal eigenvectors exhaust $(\ker K)^\perp$; Fourier expansion diagonalizes $K$.\\
\hlt{Key Insight}: Nonzero compact self-adjoint eigendirections are finite or countable, accumulating only at $0$.
\end{proof}
\end{theorem}

\begin{perspective} \label{per:hilbert-operator-spectral-theory} \hlt{Hilbert Operator Spectral Theory}\\
On a complex Hilbert space $\mathcal H$, normal $T\in\LL(\mathcal H)$ is diagonalized by spectral functional calculus.\\
By Theorem \ref{thm:hilbert-schmidt-compact-self-adjoint-spectral-theorem}, compact self-adjoint $K\in\LL(\mathcal H)$ has nonzero eigenvalues $(\lambda_k)$ and an orthonormal eigenfamily $(\psi_k)$ spanning $(\ker K)^\perp$, so $\mathcal H=\ker K\oplus\operatorname{cspan}\{\psi_k\}$ and $Kh=\sum_k\lambda_k\innerprod{h,\psi_k}\psi_k$ $\forall h\in\mathcal H$.\\
The sum is finite iff $\operatorname{rank}K<+\infty$; otherwise $\lambda_k\to0$.
\end{perspective}
